\section{はじめに}
%社会的背景
近年，教育現場においてICT(Information and Communication Technology)やEdTech(Education Technology)の活用が進んでおり，
幼児期における音楽活動は，身体的・情緒的・社会的な発達を促す重要な役割を果たすことが知られている．
茂野\cite{shigeno2020}による乳幼児のリズム同期に関する発達過程の研究によれば，3〜4歳児の段階では一定のリズムに自律的に同期することは困難であり，5〜6歳であっても
その能力は十分に発達していない．幼児は周囲の音や動きに「共振」することから表現を始め，他者の意識的な「模倣」を繰り返すステップを経て，初めて正確なリズム同期
へと発達していくとされる\cite{shigeno2020}．
このような集団でのリズム同期の獲得は，単なる音楽スキルの向上に留まらず，他者との関係性の構築や社会性の発達にも寄与することが示唆されており，
保育・教育現場における本質的な役割をも担っている．更に高橋・篠田\cite{takasino}らの研究では，幼児がリズムパターンを記憶・再生する際，演奏者の手元や人形の動き
といった「視覚的運動情報」が同時に提示されることで，再生を用意にする強い外敵援助効果(再生の容易化)をもたらすことが実証されている．したがって，音(聴覚情報)だけでなく，音楽と連動して動く人形(視覚的演出)
を組み合わせた多感覚的な自動演奏システムを導入し，子どもたちが音や動きを互いに模倣し合える環境を提供することは，幼児のリズム発達及び社会性の育成において極めて有効であると考えられる．
このような多感覚演奏システムを教育現場へ広く普及かつ実用化するためには，導入コストや保守性の観点から安価かつ省スペースな汎用マイコン(Arduino等)を複数台用いた分散配置型のシステム構成が望ましい．
しかし，多感覚的な自動演奏システムを，複数の独立したマイコンとWi-Fi通信を用いて構築する場合，ネットワーク同期において極めてシビアなリアルタイム性能が要求される．
各デバイスが独立した処理を行いながら精密にタイミングを合わせる必要があるシステムでは，わずか数ミリ秒のジッター(タイミングの揺らぎ)が致命的な音楽的破綻を招くことが先行研究でも指摘されている\cite{Lazzaro}\cite{Rottondi}．
特に不安定な無線ネットワーク環境化においては，サーバからの同期信号を基にクライアント間でいかに精密な同期を維持・確立するかが技術的課題となる．
もし，通信ジッターに拠って音と人形の動きの間，あるいはデバイス間で数十ミリ秒以上のズレが生じると，幼児は音楽と人形の動きの因果関係を正しく認識できず，模倣行動が困難となる．
これは，高橋ら\cite{takasino}の研究で示されたリズム知覚への「内的クロック」の適合を阻害し，発達途上にある子どもたちのテンポ間をかえって混乱させ，知的・教育システムとしての
有効性を損なうという問題へと直結する．
そこで本プロジェクトでは，サーバ手動の長周期同期ストリーミングを用いてネットワーク上の通信負荷とパケット送信回数を最適化
することで，Wi-Fi通金におけるジッター発生を極限まで抑制し，複数台のArduinoによるズレのない自動演奏および物理演出を実現することを目的とした．
具体的には，制約の多い低リソースなハードウェアや不安定な通信環境下であっても，サーバからのデータ転送の周期設計とクライアント側での自律的なタイミング制御を組み合わせることで，
人間のリズム知覚(特に発達段階にある幼児の知覚)では検知できないレベルの高精度な同期制御が可能であることを実証する．これにより，教育現場における実用的な多感覚演奏・協調学習システムの技術的基盤を確立することを
目指し，開発を行った．

%既存技術とプロジェクトの独自性
現在，ネットワークを介したデバイス間の同期手法として，NTP(Network Time Protocol)その軽量版であるSNTP(Simple Network Time Protocol)が広く用いられている\cite{NTPSNTP}．
しかし，これらは一般的な時刻合わせを主目的としており，無線LAN環境化では非対称な通信遅延やパケットの揺らぎによって数十ミリ秒以上の誤差が生じやすい．そのため，自動演奏に於いて要求されるシビアな同期精度を満たすことは困難である．
より高精度な同期規格として，マイクロ秒単位の同期が可能なPTP(Precision Time Protocol)\cite{PTP}があるが，これは有線LAN環境や，専用のハードウェアタイムスタンプ機能を持つ高度なネットワーク機器を前提として設計されている．
そのため，Arduinoのような低リソースなマイコンと汎用Wi-Fiモジュールを組み合わせた安価な構成においては，処理負荷の観点から実装が非現実的である．
さらに，これらの既存技術はいずれも高頻度に通信を行い，常にネットワーク経由で誤差を修正するというアプローチを取るため，無線チャネルの輻輳やマイコンの割り込み処理によるジッターの悪化を根本的に解決することができない．これに対し，本プロジェクトの最大の独自性は，同期の維持を高ヒントなネットワーク通信に依存せず，長周期のUDPブロードキャスト通信とクライアント側での自律的なタイミング補間の組み合わせによって解決する点にある．
具体的には，サーバは，頻繁な時刻同期を行うのではなく，現在のBPMに基づいて算出した1小節ごとのタイミングでのみ，小節番号を送信する．
この際，各クライアントへ個別に送信するのではなく，UDPを用いたローカルブロードキャストを採用することで，サーバからの送信時間差を完全に排除している．
更に，通信データに関して小節番号とBPM値を閾値で区別可能な1Byteのバイナリデータに圧縮することで，通信負荷を最小化し，無線LAN環境下でのパケットロスやジッターの発生を抑制している．
また，クライアント側の処理及びエラーハンドリングにおいても独自のアプローチを取る．
小節番号やBPM値を受信した各クライアントは，受信データから書く音符の長さを算出し，次回の小節同期信号を受信するまでの間，マイコンの内部タイマーを用いて自律的かつ連続的に演奏タイミングを制御する．
この際，パケットロスが生じた場合二，再送処理による遅延データの受信が致命的な楽曲のズレを引き起こすことを防ぐため，あえて再送機構を持たないUDP通信を採用している．
更に，アプリケーション層での再送要求も意図的に排除している．パケットラ筋は該当小節の演奏を行わず，次の小説信号で復帰させることで，楽曲全体のハーモニー崩壊を防ぐという，音楽的特性を活かしたフェイルセーフな設計となっている．
既存技術がネットワークに依存した密な同期を目指して破綻していたのに対し，本手法は通信を最小化し，クライアント側での自律的なタイミング制御を組み合わせることで，低リソースなマイコンと不安定な無線LAN環境下においても，幼児のリズム知覚では検知できないレベルの高精度な同期制御を実現することが可能となり，
本プロジェクトの独自性と優位性を示している．